
\documentclass{article}
\usepackage{colortbl}
\usepackage{makecell}
\usepackage{multirow}
\usepackage{supertabular}

\begin{document}

\newcounter{utterance}

\centering \large Interaction Transcript for game `cladder', experiment `full\_v1.5\_default', episode 7019 with BSU{-}SLIM\_\_Qwen3.5{-}9B{-}Base\_\_playpen{-}prm{-}9b{-}best\_of\_n.
\vspace{24pt}

{ \footnotesize  \setcounter{utterance}{1}
\setlength{\tabcolsep}{0pt}
\begin{supertabular}{c@{$\;$}|p{.15\linewidth}@{}p{.15\linewidth}p{.15\linewidth}p{.15\linewidth}p{.15\linewidth}p{.15\linewidth}}
    \# & \multicolumn{2}{c}{Player} && \multicolumn{2}{c}{Game Master} \\
    \hline

    \theutterance \stepcounter{utterance}  
    & & & \multicolumn{4}{p{0.6\linewidth}}{
        \cellcolor[rgb]{0.9,0.9,0.9}{
            \makecell[{{p{\linewidth}}}]{
                \texttt{\tiny{[P1$\langle$GM]}}
                \texttt{You are an expert in causal inference. The following question is not a typical commonsense query, but rather a meticulously designed question created by a professor specializing in causal inference, intended to assess the students' mastery of the course content.} \\
\texttt{Imagine a self{-}contained, hypothetical world with only the following conditions, and without any unmentioned factors or causal relationships: Hwax has a direct effect on gyzp and jyka. Jyka has a direct effect on lirg. Gyzp has a direct effect on lirg. Method 1: We look at how jyka correlates with lirg case by case according to gyzp. Method 2: We look directly at how jyka correlates with lirg in general. To understand how jyka affects lirg, is it more correct to use the Method 1 than Method 2?} \\
\texttt{Start your answer with 'Yes' or 'No', followed by additional reasoning or evidence to support your explanation.} \\
            }
        }
    }
     \\ \\

    \theutterance \stepcounter{utterance}  
    & \multicolumn{4}{p{0.6\linewidth}}{
        \cellcolor[rgb]{0.561,0.737,0.561}{
            \makecell[{{p{\linewidth}}}]{
                \texttt{\tiny{[P1$\rangle$GM]}}
                \texttt{No.} \\
\\ 
\texttt{In this hypothetical world, **gyzp** acts as a confounder for the relationship between **jyka** and **lirg**. Here is the causal structure:} \\
\texttt{1.  **Hwax** \$\textbackslash{}rightarrow\$ **gyzp**} \\
\texttt{2.  **Hwax** \$\textbackslash{}rightarrow\$ **jyka**} \\
\texttt{3.  **gyzp** \$\textbackslash{}rightarrow\$ **lirg**} \\
\texttt{4.  **jyka** \$\textbackslash{}rightarrow\$ **lirg**} \\
\\ 
\texttt{When you use **Method 2** (looking at the correlation between jyka and lirg in general), you are observing the total effect, which is a mixture of the direct effect (jyka \$\textbackslash{}rightarrow\$ lirg) and the indirect effect mediated through gyzp. Because gyzp is caused by Hwax, it is a common cause of both jyka and lirg. This creates a "back{-}door path" (Hwax \$\textbackslash{}rightarrow\$ jyka \$\textbackslash{}rightarrow\$ lirg and Hwax \$\textbackslash{}rightarrow\$ gyzp \$\textbackslash{}rightarrow\$ lirg) that introduces spurious correlation. Consequently, the raw correlation in Method 2 will likely be biased and will not accurately reflect the specific causal influence of jyka on lirg.} \\
\\ 
\texttt{**Method 1** attempts to control for this confounding by looking at the correlation between jyka and lirg "case by case according to gyzp." In causal inference terms, this corresponds to **stratification** or **conditioning** on the confounder (gyzp). By fixing the value of gyzp, you block the back{-}door path (Hwax \$\textbackslash{}rightarrow\$ gyzp \$\textbackslash{}rightarrow\$ lirg) because the variation in gyzp is removed. This allows the remaining observed correlation to isolate the direct effect of jyka on lirg.} \\
\\ 
\texttt{Therefore, to understand the specific causal effect of jyka on lirg and avoid the bias introduced by the common cause (gyzp), it is correct to use **Method 1**. Method 2 fails because it conflates the direct causal effect with the spurious correlation induced by the confounder.} \\
            }
        }
    }
    & & \\ \\

    \theutterance \stepcounter{utterance}  
    & & & \multicolumn{2}{p{0.3\linewidth}}{
        \cellcolor[rgb]{0.9,0.9,0.9}{
            \makecell[{{p{\linewidth}}}]{
                \texttt{\tiny{[GM$|$GM]}}
                \texttt{No} \\
            }
        }
    }
    & & \\ \\

    \theutterance \stepcounter{utterance}  
    & & & \multicolumn{2}{p{0.3\linewidth}}{
        \cellcolor[rgb]{0.9,0.9,0.9}{
            \makecell[{{p{\linewidth}}}]{
                \texttt{\tiny{[GM$|$GM]}}
                \texttt{yes} \\
            }
        }
    }
    & & \\ \\

    \theutterance \stepcounter{utterance}  
    & & & \multicolumn{2}{p{0.3\linewidth}}{
        \cellcolor[rgb]{0.9,0.9,0.9}{
            \makecell[{{p{\linewidth}}}]{
                \texttt{\tiny{[GM$|$GM]}}
                \texttt{game\_result = LOSE} \\
            }
        }
    }
    & & \\ \\

\end{supertabular}
}

\end{document}
